\documentclass[conference]{IEEEtran}

\usepackage{iftex}
\ifPDFTeX
\usepackage[utf8]{inputenc}
\usepackage[T1]{fontenc}
\fi
\usepackage{amsthm}
\usepackage{amsmath,bm,amssymb}
\usepackage{newtxtext,newtxmath}
\interdisplaylinepenalty=2500
\usepackage{xcolor}
\usepackage{graphicx}
\usepackage{pgf}
\usepackage{array,booktabs,adjustbox}
\usepackage{mathtools}
\usepackage{siunitx}
\usepackage[final]{microtype}
\usepackage{cite}
\usepackage[breaklinks=true,hidelinks]{hyperref}
\usepackage[nameinlink,capitalise,noabbrev]{cleveref}

\DeclareRobustCommand{\vct}[1]{\bm{#1}}
\DeclareRobustCommand{\mat}[1]{\bm{#1}}
\newcommand{\R}{\mathbb{R}}
\newcommand{\C}{\mathbb{C}}
\newcommand{\Hr}{^{\mathsf H}}
\newcommand{\T}{^{\mathsf T}}
\renewcommand{\Re}{\operatorname{Re}}
\renewcommand{\Im}{\operatorname{Im}}
% Keep this: PGF figures from matplotlib reference \mathdefault.
\providecommand{\mathdefault}[1]{#1}
\sisetup{detect-weight=true,detect-family=true,per-mode=symbol}
\emergencystretch=2em
\allowdisplaybreaks

\newtheoremstyle{bodysmall}{3pt}{3pt}{\normalfont\small}{0pt}{\bfseries\small}{.}{0.5em}{}
\theoremstyle{bodysmall}
\newtheorem{proposition}{Proposition}
\newtheorem{lemma}{Lemma}
\newtheorem{remark}{Remark}

% Figures are PGF drawn by plots/kalman_carma_shared_modes/draw_plots.sh from
% the committed study data (reports/results/kalman_carma_shared_modes) and
% moved here by CI; tables are the generated csm-*.tex-part files
% (tools/carma_shared_modes_study.py report).
% Every include is guarded so the article compiles from a clean checkout.
\newcommand{\ChartFigure}[2][\linewidth]{%
  \IfFileExists{#2}{\resizebox{#1}{!}{\input{#2}}}%
                   {\fbox{\parbox{0.92\linewidth}{\centering\small
                     Figure \texttt{\detokenize{#2}} not generated.\\
                     Run \texttt{plots/kalman\_carma\_shared\_modes/draw\_plots.sh}.}}}}
\newcommand{\InputPart}[1]{%
  \IfFileExists{#1}{\input{#1}}{\fbox{\small Table \texttt{\detokenize{#1}} not generated.}}}

\title{A Minimal Realization of Shared Directional CARMA Modes for the QMEKF:
Unobservable Directions, Sea-State Adaptation, Accuracy and Embedded Cost
(experimental)}
\author{%
  \IEEEauthorblockN{Mikhail S. Grushinskiy}
  \IEEEauthorblockA{Independent Researcher, USA, Fair Lawn\\
  Bareboat Necessities GitHub Project\\
  Email: mgrouch@users.github.com}%
}

\begin{document}
\maketitle

\begin{abstract}
The shared-modes QMEKF carries, on every translational axis, the real part
of one complexified CARMA(4,1) mode driven by circular complex noise with a
Hermitian excitation. Realified, the mode has eight real states per axis. We
show that exactly two directions per axis are unobservable, for every
excitation and every attitude, and that they lie in the imaginary half of the
mode. The observable quotient has six states per axis and gives a 27-state
filter. At fixed model parameters its estimates, covariances and innovations
are those of the realified filter, and amplitude, coupling and fallback
changes keep this exact; a change of the frequency scale $\kappa$ moves the
unobservable subspace by $0.45^\circ$ per percent, so each six-state
realization is its own adaptive model. We single out a cascade: the CARMA3
baseline's own coordinates driven through a two-state lossless colouring.
With the coupling off it is the baseline exactly in law, adaptation
included, and it is as well conditioned. In paired simulations on 26 seas
$\times$ 30 seeds the coupled 27-state filter lowers the 3-D position error
of the baseline on 25 seas (median $3.2\,\%$, up to $15.8\,\%$), stays
consistent, and withdraws every deliberately wrong prior; it is $58\,\%$
worse on a hull whose heave response is half the prior's, an over-coupling
of the excitation mapping. Gated on the pinned dataset, its 3-D error is
7--9\,\% below TFG-CARMA3's. On the ESP32-S3 it needs 1.9 times the
baseline's multiply-adds per sample and \SI{12.5}{kB} of filter memory. The
minimal stochastic order of the three-axis law is
$4\operatorname{rank}\Sigma+2\operatorname{rank}\Im\Sigma$, 16 motion states
(25 in all); that realization depends on the excitation, and we argue why 27
is the better operating point.
\end{abstract}

\section{Questions}
The shared-modes QMEKF \texttt{KalmanCarmaSharedModes} (CSM) is derived in
the companion article \texttt{kalman\_carma\_shared\_modes-math}, which
places the complexified mode among the complex-OU factorizations of CARMA
processes \cite{fasen21}. Its law is that of one complexified CARMA(4,1)
mode per axis; realified, that mode has eight real states. This article
answers, with proofs where they exist and with paired simulations where they
do not:
(i) Which of the eight states per axis can any measurement see?
(ii) Which six-state realization preserves the three-dimensional law, the
cross-axis co- and quadrature spectra and the Hermitian noise structure, and
what does it do under sea-state adaptation?
(iii) How does the resulting 27-state filter compare with the
independent-axis CARMA3 baseline and with TFG-CARMA3?
(iv) What does it cost on the ESP32-S3?
(v) Is 27 the right dimension?

\section{The Realified Complexified Mode}
Per axis $a\in\{x,y,z\}$ the complex kinematic state is
$\vct\zeta_a=[\zeta,\dot\zeta,\ddot\zeta,l]\T\in\C^4$,
$\dot{\vct\zeta}_a=\mat A_c\vct\zeta_a+\vct e_3\xi_a$, with $\zeta=B(s)l$,
$B(s)=s+\omega_{z0}$, and $\mat A_c$ the companion form of $P_+D$ in these
coordinates. Here $P=P_+P_-$ is the C41J quartic, $P_+$ holds its
upper-half-plane roots $\mu_1,\mu_2$, $P_-$ (conjugate coefficients) the
others, and the real Hurwitz quadratic $D$ satisfies
$|D(i\omega)|^2=\tfrac12(|P_+(i\omega)|^2+|P_-(i\omega)|^2)$. The real state
is $\vct x_a=[\Re\vct\zeta_a;\Im\vct\zeta_a]\in\R^8$ with the realified
$\mat A_8$, and the axes share it: $\mat F=\mat I_3\otimes\mat\Phi_8$,
\begin{equation}
\mat Q_{ab}=\Re(\Sigma_{ab})\mat Q_8+\Im(\Sigma_{ab})\mat J\mat Q_8,
\label{eq:Q33}
\end{equation}
$\mat J=\left[\begin{smallmatrix}\mat0&-\mat I\\\mat I&\mat0\end{smallmatrix}\right]$,
$\mat\Sigma\succeq0$ Hermitian. Every accelerometer row reads $\Re\ddot\zeta$
of the three axes through $\mat R_{wb}$. For C41J ($\omega_{z0}=0$) the roots
of $D$ are $-7.93$ and $-1.25$\,rad/s at $\kappa=1$. A QMEKF on this
realization would have 33 error states.

\section{Unobservable Directions}
\label{sec:unobs}
Write $\vct w(r)=[B(r),rB(r),r^2B(r),1]\T$; it is the eigenvector of
$\mat A_c$ for the eigenvalue $r$ (with $l=1$).
\begin{proposition}[Unobservable subspace]
\label{prop:U}
Let the roots $\mu_1,\mu_2,\bar\mu_1,\bar\mu_2$ of $P$ and $r_3,r_4$ of $D$
be distinct and $B(\mu_k)\neq0$. The unobservable subspace of
$(\mat A_8,[\vct e_3\T,\vct 0\T])$ is
\begin{equation}
\mathcal U=\{0\}\times\bigl(\operatorname{span}_{\C}\{\vct w(r_3),\vct w(r_4)\}\cap\R^4\bigr),
\end{equation}
of real dimension two. It contains no real-half component, it is
$\mat A_8$-invariant, and its annihilator in the imaginary half is spanned by
the rows $\vct n\T$ whose polynomial
$\nu(s)=n_0B+n_1sB+n_2s^2B+n_3$ is $D(s)(\alpha+\beta s)$.
\end{proposition}
\begin{proof}
A real state is a $\vct z\in\C^4$; expand $\vct z=\sum_k\gamma_k\vct w_k$ in
the eigenbasis of $\mat A_c$ (eigenvalues $\mu_1,\mu_2,r_3,r_4$). The output
is $\Re(\vct e_3\T e^{\mat A_ct}\vct z)=\tfrac12\sum_k(\gamma_kc_ke^{\lambda_kt}
+\bar\gamma_k\bar c_ke^{\bar\lambda_kt})$ with $c_k=\lambda_k^2B(\lambda_k)$.
The exponents $\mu_k,\bar\mu_k$ occur once each, so $\gamma_1=\gamma_2=0$.
Because $D$ is real, $\{r_3,r_4\}$ is closed under conjugation, and the
coefficient of $e^{r_3t}$ is $\gamma_3c_3+\bar\gamma_jc_3$ where
$r_j=\bar r_3$ ($\vct w(\bar r)=\overline{\vct w(r)}$). It vanishes iff
$\gamma_j=-\bar\gamma_3$, and then $\vct z=\gamma_3\vct w_3-\overline{\gamma_3\vct w_3}=2i\Im(\gamma_3\vct w_3)$
(for a real root $j=3$ and $\gamma_3\in i\R$). Hence $\Re\vct z=0$ and
$\Im\vct z$ ranges over the stated two-dimensional real space; conversely
such states give zero output. Invariance holds because $\mathcal U$ is a sum
of real eigenspaces. Finally $\vct n\T\vct w(r)=\nu(r)$, a real polynomial
of degree at most three that vanishes at both roots of $D$ iff $D\mid\nu$.
\end{proof}
Each root of $D$ is an eigenvalue of both $\mat A_c$ and $\bar{\mat A}_c$, so
$\mat A_8$ has it twice with a two-dimensional eigenspace, which one scalar
output cannot resolve. That is the mechanism. The rows with
$(\alpha,\beta)=(1,0),(0,1)$ read $\Im m$ and $\Im\dot m$ of the
\emph{$D$-free mode} $m=D(\tfrac d{dt})l=\xi/P_+$, so the quotient
coordinates are $[\Re\vct\zeta_a;\Im m;\Im\dot m]$. Numerically
(\texttt{tools/carma\_shared\_modes\_analysis.py}) the normalized
observability matrix has singular values down to $5\times10^{-4}$ and then
$4\times10^{-17}$, the numerical null space has no real-half component
($3\times10^{-14}$), and it agrees with the analytic $\mathcal U$ to
$3\times10^{-12}$\,degrees.

Two of the four imaginary directions are unobservable for \emph{every}
$\mat\Sigma$; the other two are observable for every $\mat\Sigma$, because
they drive $\Re\vct\zeta$ through $\Im\mat A_c$. Whether the law of the output
needs them is a stochastic question (\cref{prop:order}).

\begin{proposition}[Exactness at fixed parameters]
\label{prop:exact}
Let $\mat\Pi=\operatorname{blkdiag}(\mat I_6,\mat I_3\otimes\mat\Pi_8,\mat I_3)$
with $\ker\mat\Pi_8=\mathcal U$. For a fixed $\kappa$ and \emph{any} sequence
of excitations $\mat\Sigma_k$ (amplitude commits, estimated coupling,
coupling confidence), the 27-state QMEKF on the observable quotient, started
from $\mat\Pi\mat P_0\mat\Pi\T$, produces the same attitude, bias, $p$, $v$,
$a$ estimates, the same innovations and NIS, and the covariance
$\mat\Pi\mat P_k\mat\Pi\T$ of the QMEKF on the realified mode, including
through multiplicative resets and the coupling fallback.
\end{proposition}
\begin{proof}
In coordinates adapted to $\mathcal U$ the linear block is lower block
triangular: $\mathcal U$ is $\mat\Phi_8$-invariant (\cref{prop:U}) and
$\mat\Phi$ does not depend on $\mat\Sigma$. Every measurement row annihilates
$\mat I_3\otimes\mathcal U$: the accelerometer rows read $\Re\ddot\zeta$ through
any rotation, the magnetometer reads no motion state, and readouts of $p,v$
lie in the row space of the observability matrix (as $\mat A_8$ is
invertible, $\vct e_1\T=\vct e_3\T\mat A_8^{-2}$ on the motion block). The
Riccati recursion of the observable coordinates is then closed in
$\mat\Pi\mat P\mat\Pi\T$ with $\mat\Pi\mat Q\mat\Pi\T$, whatever $\mat Q_k$ is,
and the gains on the observable coordinates coincide \cite{kalman63,andersonmoore79}.
The EKF linearizes at the attitude, bias and acceleration estimates, which
are observable, so the linearization points agree by induction. The reset
congruence acts on the attitude rows and columns only. The fallback
$\mat P_{hz}\leftarrow\mat E+\alpha(\mat P_{hz}-\mat E)$,
$\mat E=\mat P_{hr}\mat P_{rr}^{-1}\mat P_{rz}$, maps to
$\mat\Pi_h\mat P_{hz}\mat\Pi_z\T$ and $\mat\Pi_h\mat E\mat\Pi_z\T=(\mat\Pi_h\mat P_{hr})\mat P_{rr}^{-1}(\mat P_{rz}\mat\Pi_z\T)$,
so it commutes with $\mat\Pi$. The stationary prior at Live entry is
$\mat\Pi\mat P_\infty\mat\Pi\T$ by uniqueness of the Lyapunov solution.
\end{proof}

\section{The Six-State Cascade}
The observable quotient of an axis is $\mat A_6=\mat\Pi_8\mat A_8\mat L$
(any right inverse $\mat L$, since $\mathcal U$ is invariant), input
$\mat G_6=\mat\Pi_8\mat G_8$, $\mat G_8=[\vct e_3,\vct e_7]$. $\mat J$ does not
descend to it ($\mat J\mathcal U\not\subset\mathcal U$), so in every six-state
realization the Hermitian structure is carried by two discretizations
instead of \eqref{eq:Q33}:
\begin{equation}
\mat Q_{ab}=\Re(\Sigma_{ab})\mat Q_R+\Im(\Sigma_{ab})\mat Q_I,\quad
\mat W_R=g^2\mat G\mat G\T,\ \mat W_I=g^2\mat G\mat J_2\mat G\T,
\label{eq:Q27}
\end{equation}
$\mat J_2=\left[\begin{smallmatrix}0&-1\\1&0\end{smallmatrix}\right]$, with
$\mat Q_R$ symmetric and $\mat Q_I$ antisymmetric. The block $(b,a)$ receives
$\Re\Sigma_{ab}\mat Q_R-\Im\Sigma_{ab}\mat Q_I=\mat Q_{ab}\T$, so the
three-axis $\mat Q$ is symmetric and, being the projection of
\eqref{eq:Q33}, positive semidefinite. The filter uses the following
realization of the quotient.

Write $P_+=N_r+iN_i$ with real $N_r=s^2+\Re b_1s+\Re b_0$,
$N_i=\Im b_1s+\Im b_0$. Then $P=N_r^2+N_i^2$ and, with $\xi=\xi_r+i\xi_i$,
\begin{equation}
\Re\zeta=\frac{B}{P}\,\eta,\qquad \eta=V\begin{bmatrix}\xi_r\\\xi_i\end{bmatrix},\qquad
V=\frac{[N_r,\ N_i]}{D}.
\label{eq:cascade}
\end{equation}
\begin{lemma}[Lossless colouring]
\label{lem:coinner}
$V(i\omega)V(i\omega)\Hr=1$. In input-normal coordinates (controllability
Gramian $\mat I$) its realization $(\mat A_w,\mat B_w,\mat C_w,[1,0])$
satisfies $\mat A_w+\mat A_w\T=-\mat B_w\mat B_w\T$ and
$\mat C_w=-\vct e_1\T\mat B_w\T$. Under the frequency scale,
$(\mat A_w,\mat B_w,\mat C_w)(\kappa)=(\kappa\mat A_w,\sqrt\kappa\mat B_w,\sqrt\kappa\mat C_w)(1)$.
\end{lemma}
\begin{proof}
$|N_r|^2+|N_i|^2=\tfrac12(|N_r+iN_i|^2+|N_r-iN_i|^2)=|D|^2$ on the imaginary
axis (parallelogram law). The two identities are the lossless bounded-real
equalities with $\mat X=\mat I$. $V_\kappa(s)=V_1(s/\kappa)$ gives the
scaling, and $\mat X=\mat I$ is invariant under it.
\end{proof}
The cascade block per axis is the CARMA3 baseline's C41J in its own
coordinates $[p,v,a,l]$, driven at $a$ by $\eta=\xi_r+\mat C_w\vct w$, and
the two states $\vct w$:
\begin{equation}
\mat A_6=\begin{bmatrix}\mat A_4&\vct e_3\mat C_w\\\mat0&\mat A_w\end{bmatrix},\quad
\mat G_6=\begin{bmatrix}\vct e_3&\vct0\\\multicolumn{2}{c}{\mat B_w}\end{bmatrix}.
\end{equation}
The orthogonal freedom of input-normal coordinates is fixed by making
$\mat B_w$ upper triangular with a positive diagonal (a Givens rotation),
which is the same at every $\kappa$. The construction is closed-form at
$2\times2$.

\begin{proposition}[Fixed-parameter equivalence]
\label{prop:law}
For any Hermitian $\mat\Sigma\succeq0$ and fixed $\kappa$, the realified mode,
its observable quotient and the cascade generate the same Gaussian law of
$(p_a,v_a,a_a)_{a=x,y,z}$: the same three-axis spectral density, hence the
same co-spectra, quadrature spectra, coherences and phases. The three QMEKFs
give the same estimates.
\end{proposition}
\begin{proof}
All three realize the same $1\times2$ transfer function per axis (the
quotient by construction, the cascade by \eqref{eq:cascade}) with the same
$2\times2$ noise blocks $g^2[\Re\Sigma_{ab}\mat I_2+\Im\Sigma_{ab}\mat J_2]$
per axis pair, so the spectral densities coincide. A stationary Gaussian law
is determined by its spectral density. The cascade is related to the
quotient by a similarity, so \cref{prop:exact} applies to it as well.
\end{proof}
Verified against an independent eight-state realified reference to
$3\times10^{-14}$ (transfer functions, $\kappa\in\{0.6,1,1.7\}$) and
$3\times10^{-11}$ (stationary law and lag-$h$ cross-covariance, Hermitian
$\mat\Sigma$; \texttt{carma\_shared\_modes\_realization-test}).

\section{Minimal Stochastic Order}
\label{sec:order}
\begin{proposition}[Order of the three-axis law]
\label{prop:order}
With the root conditions of \cref{prop:U} and $B(r_j)\neq0$, the minimal
state dimension of a linear Markov model generating the three-axis
displacement law is
\begin{equation}
n_{\min}=4\operatorname{rank}\mat\Sigma+2\operatorname{rank}\Im\mat\Sigma\ \le 16.
\end{equation}
\end{proposition}
\begin{proof}
With $\tilde P(s)=P(-s)$, the spectral density is
\begin{equation*}
\mat S(s)=\tfrac{g^2}{2}\tfrac{B\tilde B}{D\tilde D}
\Bigl[\tfrac{\mat\Sigma}{P_+\tilde P_-}+\tfrac{\bar{\mat\Sigma}}{P_-\tilde P_+}\Bigr].
\end{equation*}
The minimal order is the McMillan degree of the stable (positive-real)
part \cite{faurre79,lindquist15}, and for simple poles that is the sum of
the residue ranks. The residue at a root of $P_+$ is proportional to
$\mat\Sigma$, and at a root of $P_-$ to $\bar{\mat\Sigma}$. At a root $r$ of $D$,
$P_-(r)\tilde P_+(r)+P_+(r)\tilde P_-(r)=2D(r)\tilde D(r)=0$, so the bracket
is $\alpha(\mat\Sigma-\bar{\mat\Sigma})=2i\alpha\Im\mat\Sigma$. A real
antisymmetric $3\times3$ matrix has rank 0 or 2.
\end{proof}
The Hankel singular values of the positive-real part confirm the count for
a full Hermitian, a wave-like (quadrature), a real, a diagonal and a rank-one
$\mat\Sigma$ (16, 16, 12, 12, 8; \cref{fig:order}). Three consequences
follow. With in-phase (real) coupling or none, 12 motion states suffice: the
CARMA3 baseline A41 and the correlated-noise Experiment~A. Quadrature
coupling, the $90^\circ$ relation of linear waves, needs four more states, in
the plane spanned by $\operatorname{range}\Im\mat\Sigma$ (the vertical plane
of propagation). The 27-state filter is the minimal realization that keeps
the axes' shared dynamics $\mat F=\mat I_3\otimes\mat\Phi_6$ for every
$\mat\Sigma$.

A 25-state filter (16 motion states) exists. Write $\Im\mat\Sigma=\lambda(\vct e_1\vct e_2\T-\vct e_2\vct e_1\T)$
and take the Schur complement $\mat\Lambda$ of $\mat\Sigma$ on the plane
$\mat E=[\vct e_1,\vct e_2]$. Then $\eta=\vct\rho\,n+\mat E\vct\psi$ splits
into a rank-one white part and a two-dimensional coloured part $\vct\psi$
(two copies of $V$, four states). It is not a better operating point.
\begin{itemize}
\item Its coordinates depend on $\mat\Sigma$: the plane and $\mat\Lambda$
  move with every coupling update, every confidence change and every heading
  change. Every 0.1\,s commit would therefore be a re-coordinatization, and the
  exactness of \cref{prop:exact} under excitation changes is lost.
\item It is undefined at $\Im\mat\Sigma=0$, which the confidence crosses at
  every warm-up and fallback.
\item The Kronecker structure is lost.
\item It saves two states, about 14\,\% of the covariance work of the 27-state
  filter. The smallest Hankel singular values of the extra states are
  $5$--$7\times10^{-6}$ of the largest (\cref{fig:order}): small, but they
  carry the quadrature that motivates the model.
\end{itemize}

\section{Sea-State Adaptation}
\label{sec:adapt}
Every 0.1\,s the CARMA3 wrapper commits a frequency scale $\kappa$ (one for
all axes) and amplitudes $\sigma_a$. The coupling excitation and the
confidence change $\mat\Sigma$. A commit replaces $\mat F,\mat Q$ and keeps the
state, so the adaptive model, unlike the frozen one, depends on the
realization: matching spectra no longer imply matching laws.

\paragraph{Excitation changes are exact}
$\mathcal U$ and $\mat F$ do not depend on $\mat\Sigma$, so \cref{prop:exact}
covers amplitude commits, coupling updates, the confidence and the
fallback.

\paragraph{Frequency changes are not}
Every realization scales exactly with $\kappa$:
$\mat A(\kappa)=\kappa\mat S\mat A(1)\mat S^{-1}$ with
$\mat S=\operatorname{diag}(\kappa^{0,1,2,-1})$ on the complexified
coordinates, $\kappa^{0,1,2,1}$ on C41J, and $g^2\propto\kappa$. But
$\mathcal U(\kappa)=\mat S\,\mathcal U(1)$ is not $\mathcal U(1)$, because the
diagonal $\mat S$ does not preserve it (\cref{fig:U}: $0.45^\circ$ per percent
of $\kappa$, $10^\circ$ at $\pm25\,\%$). After a commit $\kappa_1\to\kappa_2$ a
filter on the realified mode would carry a component in $\mathcal U(\kappa_1)$,
with a data-dependent posterior, that leaks into the new observable
coordinates. No six-state filter carries it, and the six-state realizations
differ among themselves in how a commit re-interprets their coordinates.
Each commit rule is a different time-varying model; the cascade's has a
property none of the others has.

\paragraph{Cascade}
\begin{proposition}[Coupling off: the CARMA3 baseline, adaptation included]
\label{prop:a41}
Let $\mat\Sigma(t)=\operatorname{diag}(\sigma_a^2(t))$ and $\kappa(t)$ be
piecewise constant. At every commit, let the cascade hold $[p,v,a,l]$ and the
normalized colouring $\tilde{\vct w}_a=\vct w_a/s_a$, $s_a=\sqrt{g^2(\kappa)}\sigma_a$.
Then its law of $(p,v,a,\text{measurements})$ is the time-varying law of the
CARMA3 baseline A41, and the two QMEKFs give identical estimates.
\end{proposition}
\begin{proof}
With $\tilde{\vct w}$ driven by unit white noise, $\eta_a=s_a(t)e_a(t)$,
$e=\nu_r+\mat C_w\tilde{\vct w}$. By \cref{lem:coinner} the Gramian
equation $\dot{\mat X}=\mat A_w\mat X+\mat X\mat A_w\T+\mat B_w\mat B_w\T$
keeps $\mat X=\mat I$ for any piecewise-constant $\kappa(t)$. For $u<t$,
$E[\tilde{\vct w}(t)e(u)]=\mat\Phi_w(t,u)(\mat X\mat C_w\T+\mat B_w[1,0]\T)=0$
and $E[e(t)e(u)]=0$. So $e$ is unit white noise, and $\tilde{\vct w}(t)$ is
independent of the past of $e$. At stationarity it is also independent of
$[p,v,a,l]$, which is a functional of that past. Hence $[p,v,a,l]$ is driven
by white noise of intensity $g^2(t)\sigma_a^2(t)$, exactly A41's model and
commit rule (hold the state, replace $\mat F,\mat Q$).
\end{proof}
The test reproduces A41 through jittery $\kappa$ and amplitude commits to
$4\times10^{-13}$ relative (double precision). A realization in the
complexified coordinates, holding its imaginary half at a commit, does not
have this property: its coupling-off filter is a different adaptive model
from A41 although the frozen laws agree. With the coupling on, the cascade's
commit rule extends A41's; the experiments judge it against A41.

\section{Implementation}
\texttt{SharedModesMotion<T>}
(\texttt{src/kalman\_carma\_shared\_modes/CarmaSharedModesMotion.h})
implements the Kronecker motion interface of
\texttt{Kalman3D\_Wave\_CARMA3<\dots,Motion>} with the cascade, six states
per axis, and the short-horizon coupling monitor. The QMEKF is
\texttt{KalmanCarmaSharedModes<T>} and the wrapper
\texttt{SeaStateFusion\_CarmaSharedModes}. Attitude, biases, Joseph
updates, resets, magnetometer, coupling estimate, confidence and fallback are
the CARMA3 core's and the shared-modes filter's. Four elements are specific
to the realization:
\begin{itemize}
\item a paired discretization (\texttt{discretize\_lti\_pair},
  \texttt{StepDiscretizationPair}; the exact integrals of \cite{vanloan78}
  by series and doubling) that propagates the symmetric
  $\mat Q_R$ and the antisymmetric $\mat Q_I$ of \eqref{eq:Q27} through
  \texttt{discretize\_lti}'s scaled series, doublings and jitter series;
\item a per-axis re-coordinatization hook in the CARMA3 core
  (\texttt{take\_axis\_recoordinatization}), applied as
  $\mat S\mat P\mat S\T$ right after a commit: the cascade's transform
  scales the two colouring rows by $s_a^{\text{new}}/s_a^{\text{old}}$ (applied
  as a dense $6\times6$ product);
\item exact structural zeros ($\mat B_w$ upper triangular,
  $\mat C_w=-\mat B_w\vct e_1\T$ exactly): round-off residue in a computed
  block defeats the structural scaling of the discretization and produces
  denormals, a several times slower refresh on x86;
\item closed-form $2\times2$ algebra (Lyapunov, Cholesky, Givens) in the
  cascade's construction. Eigen's generic decompositions there cost
  \SI{44}{kB} of code on the ESP32-S3.
\end{itemize}
Conditioning favours the cascade. The condition number of the stationary
correlation matrix per axis is 147 for the realified mode and 5.5 for the
cascade, equal to A41's: the colouring states are white and uncorrelated
with $[p,v,a,l]$.

\begin{table}[!t]
\caption{Verified identities (double precision unless stated).}
\label{tab:verify}
\centering\small
\begin{adjustbox}{max width=\linewidth}
\begin{tabular}{@{}l l@{}}
\toprule
property & result \\
\midrule
unobservable dimension per axis, real-half component & 2, $3\times10^{-14}$ \\
cascade vs realified mode: transfer functions & $3\times10^{-14}$ \\
stationary 3-axis law and lag (Hermitian $\Sigma$) & $3\times10^{-11}$ \\
cascade, coupling off = A41 through adaptation & $4\times10^{-13}$ (rel.) \\
$p,v,a$ jump at a commit (core path) & 0 \\
float wrapper: $P$ symmetric, PSD & exact, min.\ corr.\ eig.\ $>0$ \\
\bottomrule
\end{tabular}
\end{adjustbox}
\end{table}

\begin{figure}[!t]
\centering
\ChartFigure{analysis_unobservable_subspace.pgf}
\caption{Motion of the unobservable subspace with the frequency scale:
principal angles between $\mathcal U(\kappa)$ and $\mathcal U(1)$.}
\label{fig:U}
\end{figure}

\begin{figure}[!t]
\centering
\ChartFigure{analysis_stochastic_order.pgf}
\caption{Normalized Hankel singular values of the positive-real part of the
three-axis displacement law for five excitations; the counts above
$10^{-7}$ match $4\operatorname{rank}\Sigma+2\operatorname{rank}\Im\Sigma$.}
\label{fig:order}
\end{figure}

\section{Numerical Experiments}
\label{sec:exp}
\paragraph{Setup}
The Monte Carlo mode of \texttt{kalman\_carma\_shared\_modes-bench} runs the
carma3\_bench engine: random-phase directional seas through the benchmark
hull's RAOs (no CARMA model in the truth), one seeded IMU and magnetometer
stream per row, and every filter of the row on the same stream. The filters
are the CARMA3 baseline with Evidence excitation (A41, 21 states) and the
27-state CSM in its deployable configuration (estimated coupling and
confidence). With the coupling off CSM \emph{is} A41 (\cref{prop:a41}), so the
paired difference isolates the directional coupling. There are 26 seas in
five groups (\cref{tab:scen}): fixed seas, coupling strength and phase
(spreading $\cos^2$ to long-crested, following, beam and head seas, heave
RAO $\times0.5$ and $\times1.5$), a gradual frequency change, rapid changes
(period, direction and height in 10\,s, direction $90^\circ$ in 60\,s, onset
and cessation), and disturbances (turns, manoeuvring, coloured and
variance-cycling accelerometer noise, $\pm25\,\%$ sample jitter, magnetic
disturbances with calibration residuals, accelerometer-bias drift). Each
sea has 30 seeds, 1200\,s at 200\,Hz, scored from 300\,s. Changes are paired
over seeds; 95\,\% intervals are bootstrap over seeds, never over the
autocorrelated samples. Nothing in the filter is fitted.

\begin{table*}[!t]
\caption{3-D position RMS of A41 and CSM, and the paired changes of CSM
against A41 (\%, positive: lower error than A41) with 95\,\% bootstrap
intervals, mean over 30 seeds per sea.}
\label{tab:scen}
\centering\small
\InputPart{csm-scenarios.tex-part}
\end{table*}

\begin{table}[!t]
\caption{Mean paired change of CSM against A41 by group (\%, positive:
better), with the range over the seas of the group. The coupling group is
dominated by heave RAO $\times0.5$ (\cref{tab:scen}).}
\label{tab:groups}
\centering\small
\begin{adjustbox}{max width=\linewidth}
\InputPart{csm-groups.tex-part}
\end{adjustbox}
\end{table}

\begin{figure*}[!t]
\centering
\ChartFigure[0.9\textwidth]{csm_paired_changes.pgf}
\caption{Paired changes of heave, horizontal and 3-D RMS of CSM against A41
with 95\,\% bootstrap intervals, per sea.}
\label{fig:paired}
\end{figure*}

\begin{figure*}[!t]
\centering
\ChartFigure[0.9\textwidth]{csm_rms3.pgf}
\caption{3-D position RMS per sea (mean over 30 seeds, log scale) of A41 and
CSM.}
\label{fig:rms3}
\end{figure*}

\paragraph{Accuracy}
The 3-D RMS of CSM is significantly lower than A41's on 25 of the 26 seas,
by $+0.1$ to $+15.8\,\%$ (median $+3.2\,\%$ over the 25), and higher on one:
heave RAO $\times0.5$ ($-58\,\%$, below). Heave is better on 24 seas (spread
JONSWAP $+5.92\,\%$ $[5.68,6.17]$, regular wave $+13.2\,\%$, swell
$+17.5\,\%$) and worse on two: heave RAO $\times0.5$ and manoeuvring
($-6.1\,\%$), where the confidence withdraws the coupling during turns and the
withdrawal itself costs heave while the 3-D error still improves
($+0.6\,\%$). Horizontal RMS is better on 24 seas, unchanged under coloured
accelerometer noise, and worse on the heave RAO $\times0.5$ sea. The
coupling gains are largest where the sea is most coherent (swell $+15.8\,\%$,
regular wave $+8.4\,\%$, PM $+5.2\,\%$ in 3-D) and smallest on the very low sea
($+0.2\,\%$) and under the coloured accelerometer noise ($+0.1\,\%$), whose
error is dominated by the shared low-frequency disturbance (818 against
818\,mm). The gradual, rapid, jitter, magnetic and bias-drift seas behave like
the spread JONSWAP ($+2.6$ to $+5.3\,\%$ 3-D).

\begin{table}[!t]
\caption{Consistency, pooled over all seas and seeds: medians of the
position (3 dof) and position-velocity (6 dof) NEES and of the accelerometer
NIS, mean 95\,\% coverage of heave, failed runs.}
\label{tab:cons}
\centering\small
\begin{adjustbox}{max width=\linewidth}
\InputPart{csm-consistency.tex-part}
\end{adjustbox}
\end{table}

\begin{figure}[!t]
\centering
\ChartFigure{csm_consistency.pgf}
\caption{Position NEES: distribution per filter (left; 3 is the nominal
mean) and per run, CSM against A41 (right).}
\label{fig:cons}
\end{figure}

\paragraph{Consistency}
The pooled median position NEES is 1.652 for CSM against 1.802 for A41, the
heave 95\,\% coverage 0.912 against 0.918, and no run of either filter
failed (1560 runs; \cref{tab:cons,fig:cons}). Both filters are
conservative, as A41 is by design (its accelerometer noise is inflated; the
median NIS is 0.09 for both). On the heave RAO $\times0.5$ sea CSM is
overconfident (median NEES 14.3 against 2.8): the over-coupling below.

\begin{figure*}[!t]
\centering
\ChartFigure[0.9\textwidth]{csm_transients.pgf}
\caption{Ensemble RMS over 10 seeds (10\,s moving window) around the sea-state
changes (shaded): error against the truth of A41 and CSM.}
\label{fig:trans}
\end{figure*}

\paragraph{Transients}
Around the frequency, period, direction and height changes and the wave
onset and cessation, CSM's error envelope stays at or below A41's
(\cref{fig:trans}). During the direction change the two coincide: the
confidence withdraws the coupling while the monitor disagrees with the
prior, and the filter is then A41. After the waves stop the residual heave
error of both decays from below a millimetre. The gradual frequency change
($+5.3\,\%$ 3-D) and the three rapid changes ($+2.6$ to $+4.6\,\%$) show no
transient penalty in the scored window.

\begin{table}[!t]
\caption{Wrong coupling priors (external, built from the truth and then
reversed, turned, phase-shifted or doubled), with and without the
confidence: 3-D RMS (mm) and median position NEES, 20 seeds; first row of
each sea: A41.}
\label{tab:robust}
\centering\small
\begin{adjustbox}{max width=\linewidth}
\InputPart{csm-robust.tex-part}
\end{adjustbox}
\end{table}

\begin{figure*}[!t]
\centering
\ChartFigure[0.9\textwidth]{csm_robustness.pgf}
\caption{Wrong coupling priors: 3-D RMS of CSM with and without confidence;
dashed: A41.}
\label{fig:robust}
\end{figure*}

\paragraph{Wrong coupling}
With the confidence on, every wrong prior (reversed, turned by $90^\circ$,
phase-shifted by $45^\circ$, doubled) is withdrawn on all three seas, and CSM
then has exactly A41's error (25.7\,mm on the spread JONSWAP; the true prior
gives 24.8\,mm; \cref{tab:robust,fig:robust}). Without the confidence the
doubled prior is the worst case on the JONSWAP seas (57.0\,mm, NEES 13.9),
followed by the reversed one (29.4\,mm); on the mixed sea the reversed and
turned priors cost 21--23\,\%. A coupling frozen at 400\,s in the keel frame
through $90^\circ$ turns gives 25.4\,mm (estimated: 25.2, A41: 25.7).

\begin{table}[!t]
\caption{Regression gates of the dataset replay (eight scored records,
default seeds, magnetometer on): worst value over the records for CSM, its
fitted bar (worst plus about 0.5\,\%, the repository's convention, gated in
\texttt{run\_tests.sh}), and TFG-CARMA3 (T41, its gated Legacy excitation;
and with Evidence excitation, as CSM uses).}
\label{tab:gates}
\centering\small
\begin{adjustbox}{max width=\linewidth}
\InputPart{csm-gates.tex-part}
\end{adjustbox}
\end{table}

\paragraph{Dataset replay and regression gates}
The eight scored records of the pinned dataset are replayed with the seeds,
sensor model and adapter configuration of every CARMA filter
(\cref{tab:dataset}). CSM is gated by its own bars (\cref{tab:gates}), and
all eight records pass. Against TFG-CARMA3 (two-frame Lie group,
independent axes) the worst heave is 19.7\,\% (JONSWAP) and 25.2\,\%
(PM-Stokes) lower. That advantage is almost entirely the excitation mode:
with Evidence excitation TFG-CARMA3 reaches 2.459 and 2.440\,\% of $H_s$,
against CSM's 2.455 and 2.403. The 3-D error is 9.4 and 7.1\,\% lower in
either mode, which is the coupling. TFG-CARMA3 is better on pitch
(4.3\,\%) and on the 3-D accelerometer and gyroscope biases (1.8 and
2.5\,\%): its two-frame Lie group separates the biases slightly better.
Against CSM with the coupling off (A41 in Evidence mode) the replay shows the
Monte Carlo pattern on every record: heave 0.2--2.8\,\% and 3-D 2--10\,\%
lower.

\paragraph{An unexpected finding: the excitation mapping over-couples}
On the sea with the heave RAO halved CSM is far worse than the baseline:
A41 is 58.0\,\% $[51.8, 63.5]$ better in 3-D, and CSM is overconfident (median
NEES 14.3 against 2.8). The six removed directions are not the cause: the
excitation mapping does not depend on the realization. The measured multiple
coherence of $z$ on $(x,y)$ is unchanged by a scaling of heave (0.87 on both
seas, from the confidence monitor). The coherence of the applied excitation,
however, rises from 0.13 to 0.38. The mapping keeps the measured regression
$\vct S_{hz}/S_{zz}$, normalizes the vertical to the prior $\sigma_z^2$ and
keeps the horizontal block at the prior. The applied coherence is therefore
$\gamma^2(S_{hh}/S_{zz})/(\sigma_h^2/\sigma_z^2)$. It grows when the hull's
heave-to-surge ratio departs from the prior's, and the PSD cap then lets the
prior explain almost all horizontal motion through heave. The consistency
test compares two data-driven regressions and cannot see it. A candidate fix
is to cap the applied coherence at the measured one,
$|b|^2\le\gamma^2\sigma_h^2\sigma_z^2$; it needs its own validation.

\begin{table*}[!t]
\caption{Dataset replay (pinned oceanography-waves-lib records, default
seeds, Evidence excitation): heave RMS in \% of $H_s$ / 3-D RMS in \% of the
reference maximum, CSM and CSM with the coupling off (A41).}
\label{tab:dataset}
\centering\small
\InputPart{csm-dataset.tex-part}
\end{table*}

\section{Embedded Efficiency}
\label{sec:cost}
\begin{table}[!t]
\caption{Analytic multiply-adds (the CARMA3 core's structured
implementation): prediction plus accelerometer update per sample, the work
of one 0.1\,s model commit (discretization, re-coordinatization), and the
total per sample at 200\,Hz.}
\label{tab:cost}
\centering\small
\begin{adjustbox}{max width=\linewidth}
\InputPart{csm-cost.tex-part}
\end{adjustbox}
\end{table}

\begin{table}[!t]
\caption{ESP32-S3 (Espressif GCC esp-14.2.0\_20241119, \texttt{-Os},
arduino-esp32 flags; object of the complete wrapper): filter object size,
code, read-only data, largest stack frame and the worst static call chain
from one update.}
\label{tab:esp}
\centering\small
\begin{adjustbox}{max width=\linewidth}
\InputPart{csm-esp32s3.tex-part}
\end{adjustbox}
\end{table}

\begin{figure}[!t]
\centering
\ChartFigure{csm_cost.pgf}
\caption{Multiply-adds per sample, covariance storage and ESP32-S3 filter
object.}
\label{fig:cost}
\end{figure}

Per sample CSM needs 7\,992 multiply-adds against A41's 4\,248, and one
0.1\,s commit 19\,440 against 4\,608 (\cref{tab:cost}). The covariance grows
from 441 to 729 floats. The paired discretization keeps $5+11+11$ series
coefficients of $6\times6$. On the ESP32-S3 the filter object is 12\,504
against 6\,136\,B, the largest stack frame 1\,888 against 1\,328\,B, and the
worst static call chain of one update is unchanged (4\,864 against
4\,832\,B; \cref{tab:esp}). The code is 33\,kB (19\,\%) larger: the paired
discretization, the closed-form construction, the two coupling spectrum
estimators and the confidence. Generic Eigen decompositions in the
construction would cost 44\,kB. Every commit applies the $6\times6$ transform
of each axis densely in the core hook (5\,832 multiply-adds); the cascade's
transform is diagonal, so a diagonal path would save almost all of that. On
the host (x86-64, one job, ranking only) the QMEKF core with 10\,Hz commits
takes 3.69\,\textmu s per sample against 2.18 for A41; with x86 flush-to-zero
2.92 against 1.80 (the float discretization series underflow into denormals,
which x86 handles in microcode). The complete wrapper takes 4.61 against
3.08\,\textmu s on average, with 99th percentiles of 13.1 against
5.7\,\textmu s: the worst case is the commit step. A scalar FPU follows the
multiply-add counts more closely than SIMD x86.

\section{Conclusions}
\paragraph{Dimension}
Two of the eight real states per axis of the complexified shared mode are
unobservable for every excitation and attitude (\cref{prop:U}). The 27-state
filter carries the six observable ones and is equivalent to the realified
model at fixed parameters, and through every excitation change: amplitude
commits, coupling estimates, the confidence and the PSD-preserving fallback
(\cref{prop:exact,prop:law}). It preserves the three-dimensional law, the co-
and quadrature spectra, the Hermitian noise structure, the $p,v,a$ estimates
and the covariance. The minimal stochastic order is
$4\operatorname{rank}\mat\Sigma+2\operatorname{rank}\Im\mat\Sigma$ (25 states),
but that realization depends on the excitation and is not a better operating
point; 21 states suffice only without quadrature coupling.

\paragraph{Realization}
The cascade (the baseline's own coordinates driven through a lossless
two-state colouring) is the filter's realization: with the coupling off it is
the CARMA3 baseline exactly, adaptation included (\cref{prop:a41}), so the
fallback model is the baseline, and it is as well conditioned as the
baseline.

\paragraph{Accuracy and cost}
Against the baseline, with the same excitation mode, the coupling lowers the
3-D error on 25 of 26 simulated seas (median $3.2\,\%$, up to $15.8\,\%$),
keeps the filter consistent and withdraws every deliberately wrong prior. It
costs 1.9 times the multiply-adds per sample, twice the filter memory and
19\,\% more code on the ESP32-S3.

\paragraph{Next experiment}
The excitation mapping over-couples when the hull's heave-to-surge ratio
departs from the prior ($58\,\%$ worse than the baseline with the heave RAO
halved). Capping the applied coherence at the measured one is the next
falsifiable experiment.

\begin{thebibliography}{9}
\bibitem{andersonmoore79}
B. D. O. Anderson and J. B. Moore, \emph{Optimal Filtering}. Prentice-Hall,
1979.
\bibitem{faurre79}
P. Faurre, M. Clerget and F. Germain, \emph{Op\'erateurs rationnels
positifs}. Dunod, 1979.
\bibitem{lindquist15}
A. Lindquist and G. Picci, \emph{Linear Stochastic Systems: A Geometric
Approach to Modeling, Estimation and Identification}. Springer, 2015.
\bibitem{kalman63}
R. E. Kalman, ``Mathematical description of linear dynamical systems,''
\emph{J. SIAM Control}, vol. 1, no. 2, pp. 152--192, 1963.
\bibitem{fasen21}
V. Fasen-Hartmann and M. Scholz, ``Factorization and discrete-time
representation of multivariate CARMA processes,'' arXiv:2102.11681, 2021.
\bibitem{vanloan78}
C. F. Van Loan, ``Computing integrals involving the matrix exponential,''
\emph{IEEE Trans. Autom. Control}, vol. 23, no. 3, pp. 395--404, 1978.
\end{thebibliography}
\end{document}
